% ==============================================================================
% TRƯỜNG ĐẠI HỌC KHOA HỌC TỰ NHIÊN, ĐHQG-HCM
% KHOA CÔNG NGHỆ THÔNG TIN
% CHƯƠNG 3: NEGCL - NOISE-ENHANCED GRAPH CONTRASTIVE LEARNING
% Toàn văn hoàn chỉnh: Kiến trúc N1-N4, Phương pháp AC-DNS, Đối sánh Paper Table 4,
% Phân tích Homophily, Tài nguyên tính toán và Hướng phát triển tương lai
% ==============================================================================

\section{NEGCL: Noise-Enhanced Graph Contrastive Learning}
\label{chap:negcl}

NEGCL kết hợp biểu diễn cấu trúc tương tác, tăng cường đa phương thức bằng nhiễu ngẫu nhiên và học đối chiếu trên siêu đồ thị (Hypergraph) để xây dựng không gian biểu diễn user--item cuối cùng~\cite{shi2025negcl}. Thay vì tiếp tục sử dụng sáu cấu hình thử nghiệm cũ còn rời rạc, chương này tập trung vào bốn hướng cải tiến có mục tiêu toán học và kỹ thuật rõ ràng (N1--N4): làm sạch đồ thị KNN, điều tiết nhiễu theo node và phương thức, nhận diện false negative trong hàm mất mát, và tinh chỉnh siêu đồ thị kép thưa. 

Đồng thời, trên cơ sở phân tích đặc thù cấu trúc và độ đồng nhất đa phương thức của từng tập dữ liệu, chương này đề xuất phương pháp thống nhất \textbf{AC-DNS} (\textit{Adaptive Confidence-aware Densification and Dual Adaptive Noise Scheduling}). Phương pháp AC-DNS đã thiết lập kỷ lục hiệu năng mới, vượt qua toàn diện kết quả công bố trong bài báo gốc NEGCL (Table~4 KBS 2025~\cite{shi2025negcl}) trên cả ba tập dữ liệu chuẩn Amazon (Baby, Sports, Clothing).

\begin{table}[H]
\centering
\small
\setlength{\tabcolsep}{5pt}
\caption{Tổng quan bốn hướng cải tiến khảo sát (N1--N4) và Phương pháp Thống nhất AC-DNS.}
\label{tab:negcl_methods_overview}
\begin{tabularx}{\linewidth}{%
    >{\centering\arraybackslash}p{1.4cm}
    >{\raggedright\arraybackslash}p{4.2cm}
    X
    >{\centering\arraybackslash}p{2.2cm}}
\toprule
\textbf{Ký hiệu} & \textbf{Tên phương án} & \textbf{Vùng can thiệp trong kiến trúc} & \textbf{Tham số bổ sung} \\
\midrule
\textbf{N1} & Confidence-aware Graph Densification and Pruning & Đồ thị item--item trước khi lan truyền MGE & Không (0 tham số) \\
\textbf{N2} & Warm-up Degree--Modality Noise Router & Biên độ và lịch trình nhiễu trong MGE & 114 tham số \\
\textbf{N3} & False-Negative-Aware Ranking \& CL & Hàm mục tiêu BPR loss và Contrastive loss & Không (0 tham số) \\
\textbf{N4} & Sparse Dual-Hypergraph Refinement & Biểu diễn User/Item trước phép Fusion cuối & 2 scalar ($\gamma_u, \gamma_i$) \\
\midrule
\textbf{AC-DNS} & \textbf{Adaptive Confidence Densification + Dual Adaptive Noise} & \textbf{Hợp nhất làm dày đồ thị tin cậy thích ứng bậc node và điều chế nhiễu kép} & \textbf{Không (0 tham số)} \\
\bottomrule
\end{tabularx}
\end{table}

\begin{figure}[H]
\centering
\begin{tikzpicture}[scale=0.90, transform shape,
    block/.style={rectangle, draw, rounded corners=2pt, minimum height=0.65cm, minimum width=2.4cm, font=\scriptsize, align=center, fill=blue!8},
    newmod/.style={rectangle, draw=red!70!black, dashed, thick, rounded corners=2pt, minimum height=0.7cm, minimum width=3.0cm, font=\scriptsize\bfseries, align=center, fill=orange!15},
    propmod/.style={rectangle, draw=teal!80!black, thick, rounded corners=2pt, minimum height=0.75cm, minimum width=3.4cm, font=\scriptsize\bfseries, align=center, fill=teal!15},
    merge/.style={rectangle, draw, rounded corners=2pt, minimum height=0.7cm, minimum width=5.4cm, font=\scriptsize, align=center, fill=blue!8},
    arr/.style={-{Stealth[length=2mm]}, thick},
    modarr/.style={-{Stealth[length=2mm]}, thick, red!70!black, dashed},
]
\node[block] (interaction) at (-3.8,0) {Đồ thị tương tác\\CGE / GTLayer};
\node[block] (features) at (3.8,0) {Đặc trưng Visual / Text};
\node[newmod] (n1) at (3.8,-1.35) {[N1] Confidence-aware\\KNN + Pruning};
\node[newmod] (n2) at (3.8,-2.75) {[N2] Degree--Modality\\Noise Router};
\node[block] (mge) at (3.8,-4.15) {Multimodal Graph\\Enhancement (MGE)};
\node[newmod] (n4) at (-3.8,-4.15) {[N4] Sparse Dual\\Hypergraph $\mathcal{H}_u, \mathcal{H}_i$};
\node[merge] (fusion) at (0,-5.65) {Fusion biểu diễn cấu trúc, đa phương thức và hypergraph};
\node[newmod] (n3) at (0,-7.05) {[N3] False-Negative-Aware\\BPR + Contrastive Loss};

\draw[arr] (features) -- (n1);
\draw[arr] (n1) -- (n2);
\draw[arr] (n2) -- (mge);
\draw[arr] (interaction) -- (n4);
\draw[arr] (interaction.south) -- ++(0,-1.15cm) -| ([xshift=-1.6cm]fusion.north);
\draw[arr] (n4) -- ([xshift=-0.7cm]fusion.north);
\draw[arr] (mge) -- ([xshift=1.2cm]fusion.north);
\draw[arr] (fusion) -- (n3);
\draw[modarr] (n1.west) -- ++(-0.8cm,0) |- (n4.east);
\end{tikzpicture}
\caption{Vị trí can thiệp của bốn hướng cải tiến (N1--N4) trong luồng kiến trúc NEGCL.}
\label{fig:negcl_flow_overview}
\end{figure}

% ==============================================================================
% 3.1 CÁC HƯỚNG CẢI TIẾN ĐỀ XUẤT
% ==============================================================================
\subsection{Kiến trúc và các phương án cải tiến}
\label{sec:negcl_methods}

\subsubsection{Cải tiến 1: Confidence-aware Graph Densification and Pruning (NEGCL-N1)}

\paragraph{Động lực kỹ thuật:}
Đồ thị KNN xây dựng thuần túy từ khoảng cách cosine của đặc trưng hình ảnh và văn bản thường đưa vào nhiều cạnh kết nối giả tạo: các item có cùng tông màu hoặc cùng từ khóa mô tả chung chung nhưng thuộc các danh mục hoàn toàn khác biệt và không bao giờ được mua cùng. N1 kết hợp tín hiệu tương đồng đa phương thức với tín hiệu tương quan cộng tác lịch sử để lọc sạch đồ thị, chỉ giữ các cạnh có độ tin cậy cao và áp dụng cắt tỉa phân vị~\cite{qi2025seeing}.

\paragraph{Cơ sở toán học:}
Với hai item $i$ và $j$, độ tin cậy cạnh $c_{ij}$ được xác định bởi hàm kết hợp lồi:
\begin{equation}
c_{ij} = \eta \, s_{ij}^{\text{mm}} + (1 - \eta) \, s_{ij}^{\text{cf}}, \quad \eta \in [0, 1],
\label{eq:n1_conf_score}
\end{equation}
trong đó $s_{ij}^{\text{mm}} = \omega_v \cos(\mathbf{v}_i, \mathbf{v}_j) + \omega_t \cos(\mathbf{t}_i, \mathbf{t}_j)$ là cosine similarity đa phương thức, còn $s_{ij}^{\text{cf}} = \frac{\mathbf{R}_{:,i}^\top \mathbf{R}_{:,j}}{\|\mathbf{R}_{:,i}\| \|\mathbf{R}_{:,j}\|}$ đo mức độ tương đồng từ hành vi đồng tiêu thụ trong ma trận tương tác $\mathbf{R}$. Cạnh chỉ được giữ lại khi nằm trong Top-$K$ láng giềng và vượt qua ngưỡng phân vị cắt tỉa $\tau$:
\begin{equation}
\mathbf{A}_{ij}^{\text{conf}} = 
\begin{cases}
c_{ij}, & j \in \text{TopK}_i(c) \ \text{và} \ c_{ij} \ge \tau, \\
0, & \text{ngược lại.}
\end{cases}
\label{eq:n1_pruning}
\end{equation}
Sau khi chuẩn hóa đối xứng Laplacian $\tilde{\mathbf{A}}_{\text{conf}} = \mathbf{D}_{\text{conf}}^{-1/2} \mathbf{A}^{\text{conf}} \mathbf{D}_{\text{conf}}^{-1/2}$, item embedding được làm giàu theo cơ chế residual:
\begin{equation}
\mathbf{E}_i' = \mathbf{E}_i + \lambda_{\text{dense}} \, \tilde{\mathbf{A}}_{\text{conf}} \mathbf{E}_i.
\label{eq:n1_residual}
\end{equation}

\begin{figure}[H]
\centering
\begin{tikzpicture}[scale=0.95, transform shape,
    block/.style={rectangle, draw, rounded corners=2pt, minimum height=0.65cm, minimum width=2.7cm, font=\scriptsize, align=center, fill=blue!8},
    newmod/.style={rectangle, draw=red!70!black, dashed, thick, rounded corners=2pt, minimum height=0.7cm, minimum width=3.5cm, font=\scriptsize\bfseries, align=center, fill=orange!15},
    arr/.style={-{Stealth[length=2mm]}, thick},
]
\node[block] (mm) at (-3.4,0) {Visual/Text similarity\\$s_{ij}^{\text{mm}}$};
\node[block] (cf) at (3.4,0) {Collaborative similarity\\$s_{ij}^{\text{cf}}$};
\node[newmod] (confidence) at (0,-1.35) {[N1] Confidence score\\$c_{ij} = \eta s_{ij}^{\text{mm}} + (1-\eta)s_{ij}^{\text{cf}}$};
\node[newmod] (prune) at (0,-2.75) {Top-$K$ + threshold pruning\\$\tilde{\mathbf{A}}_{\text{conf}}$};
\node[block] (residual) at (0,-4.15) {Residual update\\$\mathbf{E}_i' = \mathbf{E}_i + \lambda_d \tilde{\mathbf{A}}_{\text{conf}} \mathbf{E}_i$};
\draw[arr] (mm) -- (confidence);
\draw[arr] (cf) -- (confidence);
\draw[arr] (confidence) -- (prune);
\draw[arr] (prune) -- (residual);
\end{tikzpicture}
\caption{Cơ chế N1: Kết hợp độ tương đồng đa phương thức và cộng tác, sau đó lọc cạnh tin cậy cao trước khi làm giàu biểu diễn item.}
\label{fig:negcl_n1}
\end{figure}

\paragraph{Chi phí tính toán:}
N1 không bổ sung bất kỳ tham số học gradient nào ($0$ tham số). Ma trận $\tilde{\mathbf{A}}_{\text{conf}}$ được tiền tính toán (pre-computed) trước khi huấn luyện và lưu trữ dưới dạng ma trận thưa (sparse tensor), giúp chi phí mỗi epoch chỉ là phép nhân ma trận thưa $\mathcal{O}(|\mathcal{E}_{\text{conf}}| \cdot d)$.

% ------------------------------------------------------------------------------
\subsubsection{Cải tiến 2: Warm-up Degree--Modality Noise Router (NEGCL-N2)}

\paragraph{Động lực kỹ thuật:}
NEGCL nguyên bản áp dụng một biên độ nhiễu tĩnh $\varepsilon$ cố định cho mọi node trong suốt quá trình huấn luyện. Tuy nhiên, một biên độ tĩnh sẽ gây xáo trộn nghiêm trọng các biểu diễn ở những epoch đầu khi mạng chưa học được cấu trúc cơ bản. Đồng thời, các node phổ biến (head nodes) đã có tín hiệu tương tác dày đặc cần được bảo tồn cấu trúc, trong khi các node đuôi dài (tail nodes) rất cần sự khuếch đại nhiễu để khám phá không gian tiềm ẩn~\cite{liu2025gated}.

\paragraph{Cơ sở toán học:}
Vector trạng thái đầu vào của node $i$ bao gồm bậc tương tác đồ thị và chuẩn độ dài các vector đa phương thức:
\begin{equation}
\mathbf{q}_i = \left[ \log(d_i + 1) \,\|\, \cos(\mathbf{e}_i^v, \mathbf{e}_i^t) \,\|\, \|\mathbf{e}_i^v\|_2 \,\|\, \|\mathbf{e}_i^t\|_2 \right].
\label{eq:n2_state_vector}
\end{equation}
Mạng Router (MLP hai tầng nhỏ gọn) sinh ra hai hệ số biên độ nhiễu bị chặn cho Visual và Text:
\begin{equation}
[\varepsilon_{i,v}, \varepsilon_{i,t}] = \varepsilon_{\text{max}} \cdot \sigma\left( \mathbf{W}_2 \text{ReLU}(\mathbf{W}_1 \mathbf{q}_i + \mathbf{b}_1) + \mathbf{b}_2 \right).
\label{eq:n2_mlp_router}
\end{equation}
Hệ số điều hòa theo thời gian Warm-up $r_t = \min(1.0, t / T_w)$ kiểm soát mức độ tác động của nhiễu tại epoch $t$:
\begin{equation}
\tilde{\mathbf{e}}_i^m = \mathbf{e}_i^m + r_t \, \varepsilon_{i,m} \, \text{sign}(\mathbf{e}_i^m) \odot \boldsymbol{\delta}_i^m, \quad \boldsymbol{\delta}_i^m \sim \mathcal{N}(\mathbf{0}, \mathbf{I}), \quad m \in \{v, t\}.
\label{eq:n2_noise_applied}
\end{equation}

\begin{figure}[H]
\centering
\begin{tikzpicture}[scale=0.95, transform shape,
    block/.style={rectangle, draw, rounded corners=2pt, minimum height=0.65cm, minimum width=2.5cm, font=\scriptsize, align=center, fill=blue!8},
    newmod/.style={rectangle, draw=red!70!black, dashed, thick, rounded corners=2pt, minimum height=0.7cm, minimum width=4.0cm, font=\scriptsize\bfseries, align=center, fill=orange!15},
    arr/.style={-{Stealth[length=2mm]}, thick},
]
\node[block] (stats) at (0,0) {Degree + Visual/Text statistics\\$\mathbf{q}_i$};
\node[newmod] (router) at (0,-1.35) {[N2] MLP Router\\$\mathbf{q}_i \to (\varepsilon_{i,v}, \varepsilon_{i,t})$};
\node[block] (warmup) at (4.2,-1.35) {Warm-up schedule\\$r_t = \min(1, t/T_w)$};
\node[newmod] (noise) at (0,-2.8) {Modality-specific bounded noise\\$r_t \varepsilon_{i,m} \boldsymbol{\delta}_i^m$};
\node[block] (mge) at (0,-4.15) {MGE Visual/Text};
\draw[arr] (stats) -- (router);
\draw[arr] (router) -- (noise);
\draw[arr] (warmup) -- (noise);
\draw[arr] (noise) -- (mge);
\end{tikzpicture}
\caption{Cơ chế N2: Router học biên độ nhiễu theo từng node và phương thức kết hợp lịch trình Warm-up.}
\label{fig:negcl_n2}
\end{figure}

\paragraph{Chi phí bổ sung:}
Với kiến trúc MLP $4 \to 16 \to 2$, N2 chỉ bổ sung đúng $4 \times 16 + 16 + 16 \times 2 + 2 = \mathbf{114}$ tham số, chiếm dung lượng bộ nhớ không đáng kể ($<1$\,KB).

% ------------------------------------------------------------------------------
\subsubsection{Cải tiến 3: False-Negative-Aware Ranking and Contrastive Learning (NEGCL-N3)}

\paragraph{Động lực kỹ thuật:}
Trong phản hồi ngầm định (implicit feedback), các mẫu chưa được người dùng tương tác được lấy mẫu ngẫu nhiên làm negative item. Tuy nhiên, nhiều item negative thực chất là các mặt hàng rất phù hợp với sở thích của người dùng nhưng họ chưa có cơ hội tiếp cận (false negative). Việc ép mô hình đẩy xa các item này trong BPR loss và Contrastive loss InfoNCE tạo ra các gradient đối kháng nghiêm trọng, phá vỡ cấu trúc cụm sở thích tiềm năng~\cite{balmaseda2025discovering}.

\paragraph{Cơ sở toán học:}
Với cặp $(p, n)$ gồm positive item $p$ và negative item $n$, xác suất false negative $q_{pn}$ được ước lượng thông qua độ tương đồng đa phương thức và độ tương đồng cấu trúc láng giềng Jaccard:
\begin{equation}
q_{pn} = \sigma\left( a \, s_{pn}^{\text{mm}} + b \, s_{pn}^{\text{nbr}} - c \right), \quad s_{pn}^{\text{nbr}} = \frac{|\mathcal{N}(p) \cap \mathcal{N}(n)|}{|\mathcal{N}(p) \cup \mathcal{N}(n)|}.
\label{eq:n3_fn_score}
\end{equation}
Trọng số của mẫu trong BPR loss được điều chỉnh giảm mềm với cận dưới bảo vệ $w_{\text{min}}$:
\begin{equation}
w_{pn} = \max(w_{\text{min}}, 1 - q_{pn}), \quad \mathcal{L}_{\text{BPR}}^{\text{FN}} = -\sum_{(u,p,n)} w_{pn} \log \sigma(\hat{y}_{up} - \hat{y}_{un}).
\label{eq:n3_bpr_loss}
\end{equation}
Tương tự, trong Contrastive Loss InfoNCE, trọng số của các negative view được giảm trừ:
\begin{equation}
\mathcal{L}_{\text{CL}}^{\text{FN}} = -\log \frac{\exp(\text{sim}(\mathbf{z}_i, \mathbf{z}_i^+) / \tau)}{\exp(\text{sim}(\mathbf{z}_i, \mathbf{z}_i^+) / \tau) + \sum_{j \ne i} (1 - q_{ij}) \exp(\text{sim}(\mathbf{z}_i, \mathbf{z}_j) / \tau)}.
\label{eq:n3_cl_loss}
\end{equation}

\begin{figure}[H]
\centering
\begin{tikzpicture}[scale=0.95, transform shape,
    block/.style={rectangle, draw, rounded corners=2pt, minimum height=0.65cm, minimum width=2.8cm, font=\scriptsize, align=center, fill=blue!8},
    newmod/.style={rectangle, draw=red!70!black, dashed, thick, rounded corners=2pt, minimum height=0.7cm, minimum width=4.0cm, font=\scriptsize\bfseries, align=center, fill=orange!15},
    arr/.style={-{Stealth[length=2mm]}, thick},
]
\node[block] (triplet) at (-3.5,0) {BPR triplet $(u,p,n)$};
\node[block] (views) at (3.5,0) {Contrastive views $(\mathbf{z}_i, \mathbf{z}_j)$};
\node[newmod] (fn) at (0,-1.45) {[N3] False-negative score\\$q = \sigma(a s^{\text{mm}} + b s^{\text{nbr}} - c)$};
\node[block] (bpr) at (-2.5,-3.0) {$\mathcal{L}_{\text{BPR}}^{\text{FN}}$};
\node[block] (cl) at (2.5,-3.0) {$\mathcal{L}_{\text{CL}}^{\text{FN}}$};
\node[newmod] (total) at (0,-4.35) {$\mathcal{L} = \mathcal{L}_{\text{BPR}}^{\text{FN}} + \lambda_c \mathcal{L}_{\text{CL}}^{\text{FN}} + \lambda_r \|\Theta\|_2^2$};
\draw[arr] (triplet) -- (fn);
\draw[arr] (views) -- (fn);
\draw[arr] (fn) -- (bpr);
\draw[arr] (fn) -- (cl);
\draw[arr] (bpr) -- (total);
\draw[arr] (cl) -- (total);
\end{tikzpicture}
\caption{Cơ chế N3: Ước lượng xác suất false negative để hiệu chỉnh trọng số BPR và Contrastive loss.}
\label{fig:negcl_n3}
\end{figure}

\paragraph{Chi phí bổ sung:}
N3 không thêm tham số học ($0$ tham số) và tận dụng lại ma trận khoảng cách tương đồng có sẵn từ N1.

% ------------------------------------------------------------------------------
\subsubsection{Cải tiến 4: Sparse Dual-Hypergraph Refinement (NEGCL-N4)}

\paragraph{Động lực kỹ thuật:}
Đồ thị tương tác thông thường chỉ mô tả quan hệ nhị phân từng cặp (pairwise). N4 thiết lập đồng thời hai siêu đồ thị thưa: User Hypergraph $\mathcal{H}_u$ gom các nhóm người dùng cùng chung sở thích đa phương thức, và Item Hypergraph $\mathcal{H}_i$ gom các nhóm sản phẩm có quan hệ tương hỗ bậc cao~\cite{guo2025mmhcl}.

\paragraph{Cơ sở toán học:}
Với ma trận liên thuộc (incidence matrix) $\mathbf{H}_u$ và $\mathbf{H}_i$, toán tử lan truyền siêu đồ thị chuẩn hóa được tính như sau:
\begin{equation}
\mathcal{P}(\mathbf{H}, \mathbf{E}) = \mathbf{D}_v^{-1/2} \mathbf{H} \mathbf{W} \mathbf{D}_e^{-1} \mathbf{H}^\top \mathbf{D}_v^{-1/2} \mathbf{E}.
\label{eq:n4_hyper_prop}
\end{equation}
Biểu diễn của User và Item được cập nhật qua liên kết thặng dư có chặn:
\begin{equation}
\mathbf{E}_u' = \mathbf{E}_u + \gamma_u \mathcal{P}(\mathbf{H}_u, \mathbf{E}_u), \quad
\mathbf{E}_i' = \mathbf{E}_i + \gamma_i \mathcal{P}(\mathbf{H}_i, \mathbf{E}_i), \quad \gamma_u, \gamma_i \in [0, 1].
\label{eq:n4_hyper_residual}
\end{equation}

\begin{figure}[H]
\centering
\begin{tikzpicture}[scale=0.95, transform shape,
    block/.style={rectangle, draw, rounded corners=2pt, minimum height=0.65cm, minimum width=2.6cm, font=\scriptsize, align=center, fill=blue!8},
    newmod/.style={rectangle, draw=red!70!black, dashed, thick, rounded corners=2pt, minimum height=0.7cm, minimum width=3.3cm, font=\scriptsize\bfseries, align=center, fill=orange!15},
    merge/.style={rectangle, draw, rounded corners=2pt, minimum height=0.7cm, minimum width=5.2cm, font=\scriptsize, align=center, fill=blue!8},
    arr/.style={-{Stealth[length=2mm]}, thick},
]
\node[block] (eu) at (-3.3,0) {User embeddings $\mathbf{E}_u$};
\node[block] (ei) at (3.3,0) {Item embeddings $\mathbf{E}_i$};
\node[newmod] (hu) at (-3.3,-1.45) {[N4] Sparse User\\Hypergraph $\mathbf{H}_u$};
\node[newmod] (hi) at (3.3,-1.45) {[N4] Sparse Item\\Hypergraph $\mathbf{H}_i$};
\node[block] (eur) at (-3.3,-2.9) {$\mathbf{E}_u' = \mathbf{E}_u + \gamma_u \mathcal{P}_u$};
\node[block] (eir) at (3.3,-2.9) {$\mathbf{E}_i' = \mathbf{E}_i + \gamma_i \mathcal{P}_i$};
\node[merge] (fusion) at (0,-4.35) {Fusion cuối và tính điểm tương tác user--item};
\draw[arr] (eu) -- (hu);
\draw[arr] (ei) -- (hi);
\draw[arr] (hu) -- (eur);
\draw[arr] (hi) -- (eir);
\draw[arr] (eur) -- (fusion);
\draw[arr] (eir) -- (fusion);
\end{tikzpicture}
\caption{Cơ chế N4: Tinh chỉnh độc lập biểu diễn User và Item qua hai siêu đồ thị thưa.}
\label{fig:negcl_n4}
\end{figure}

% ------------------------------------------------------------------------------
\subsubsection{Phương pháp Đề xuất Thống nhất: AC-DNS (Adaptive Confidence-aware Densification and Dual Adaptive Noise Scheduling)}
\label{sec:ac_dns_formulation}

\paragraph{Động lực thiết kế thống nhất:}
Từ kết quả thực nghiệm của các hướng độc lập, N1 mang lại nền tảng cấu trúc đồ thị sạch và vững chắc, trong khi N2 mang lại khả năng tạo góc nhìn đối chiếu giàu tính biểu diễn. Tuy nhiên, việc áp dụng router MLP của N2 đôi khi thiếu sự kiểm soát chặt chẽ đối với các tập dữ liệu có hành vi tiêu dùng phân mảnh (như Baby). 

Do đó, nhóm đề xuất phương pháp thống nhất \textbf{AC-DNS}, tích hợp hai cơ chế cốt lõi không tham số học:
\begin{enumerate}[leftmargin=*, label=(\arabic*)]
  \item \textbf{Làm dày đồ thị tin cậy thích ứng theo bậc node (Degree-Aware Adaptive Confidence Densification):} Hệ số tin cậy tương quan cộng tác $\gamma_{\text{cf}}(i)$ được tính thích ứng dựa trên bậc tương tác chuẩn hóa của từng node:
  \begin{equation}
    \gamma_{\text{cf}}(i) = \gamma_{\text{base}} \cdot (0.5 + \text{norm\_deg}_i), \quad \text{norm\_deg}_i = \frac{d_i - d_{\text{min}}}{d_{\text{max}} - d_{\text{min}} + \epsilon}.
    \label{eq:ac_dns_gamma_deg}
  \end{equation}
  Độ tin cậy cạnh giữa item $i$ và $j$ được xác lập:
  \begin{equation}
    c_{ij}^{\text{adapt}} = (1 - \bar{\gamma}_{ij}) \, s_{ij}^{\text{mm}} + \bar{\gamma}_{ij} \, s_{ij}^{\text{cf}}, \quad \bar{\gamma}_{ij} = \frac{\gamma_{\text{cf}}(i) + \gamma_{\text{cf}}(j)}{2}.
  \end{equation}
  Ma trận kề làm dày $\tilde{\mathbf{A}}_{\text{dense}}$ được chuẩn hóa Laplacian và cập nhật vào biểu diễn item theo trọng số $\lambda_{\text{dense}}$.

  \item \textbf{Điều chế nhiễu thích ứng kép (Dual Adaptive Noise Scheduling):} Kết hợp đồng thời điều chế theo \textit{Trục thời gian} (Warm-up Schedule $r_t$) và theo \textit{Không gian cấu trúc} (Degree Scaling Factor $s_{\text{deg}}(i)$):
  \begin{equation}
    r_t = \min\left(1.0, \frac{t}{T_{\text{warmup}}}\right), \quad 
    s_{\text{deg}}(i) = 1.2 - 0.4 \cdot \text{norm\_deg}_i \in [0.8, 1.2].
    \label{eq:ac_dns_dual_noise}
  \end{equation}
  Biên độ nhiễu thực tế tại epoch $t$ cho node $i$ và phương thức $m \in \{v, t\}$ là:
  \begin{equation}
    \varepsilon_i^m(t) = r_t \cdot \varepsilon_{\text{base}}^m \cdot s_{\text{deg}}(i).
  \end{equation}
  Cơ chế này đảm bảo: các node phổ biến (Head items) nhận mức nhiễu bảo vệ nhỏ ($0.8 \times \varepsilon_{\text{base}}$), trong khi các node đuôi dài (Tail items) nhận mức nhiễu khuếch đại ($1.2 \times \varepsilon_{\text{base}}$) để kích thích khám phá không gian tiềm ẩn.
\end{enumerate}

\begin{figure}[H]
\centering
\begin{tikzpicture}[scale=0.92, transform shape,
    block/.style={rectangle, draw, rounded corners=2pt, minimum height=0.75cm, font=\scriptsize, align=center, fill=blue!8},
    propmod/.style={rectangle, draw=teal!80!black, thick, rounded corners=2pt, minimum height=1.15cm, text width=5.6cm, inner sep=5pt, font=\scriptsize\bfseries, align=center, fill=teal!15},
    arr/.style={-{Stealth[length=2mm]}, thick},
]
\node[block, text width=8cm] (inputs) at (0,0)
{Bậc tương tác $d_i$ và Visual/Text Features};
\node[propmod] (dens) at (-3.5,-1.9)
{[AC-DNS Part 1]\\Adaptive Confidence Graph\\
$\gamma_{\text{cf}}(i) = \gamma_{\text{base}}(0.5 + \text{norm\_deg}_i)$};
\node[propmod] (noise) at (3.5,-1.9)
{[AC-DNS Part 2]\\Dual Adaptive Noise\\
$\varepsilon_i(t) = r_t \varepsilon_{\text{base}} s_{\text{deg}}(i)$};
\node[block, text width=7.4cm] (mge) at (0,-3.8)
{Multimodal Graph Enhancement (MGE) \& GTLayer};
\node[block, text width=9cm] (loss) at (0,-5.1)
{Fusion Representation\\Joint BPR + Hypergraph Contrastive Loss};

\draw[arr] ([xshift=-3.5cm]inputs.south) -- (dens.north);
\draw[arr] ([xshift=3.5cm]inputs.south) -- (noise.north);
\draw[arr] (dens.south) -- ++(0,-0.45) -| ([xshift=-2.2cm]mge.north);
\draw[arr] (noise.south) -- ++(0,-0.45) -| ([xshift=2.2cm]mge.north);
\draw[arr] (mge.south) -- (loss.north);
\end{tikzpicture}
\caption{Kiến trúc phương pháp đề xuất AC-DNS: Hợp nhất làm dày đồ thị tin cậy theo bậc node và điều chế nhiễu kép.}
\label{fig:negcl_ac_dns}
\end{figure}

% ==============================================================================
% 3.2 KẾT QUẢ THỰC NGHIỆM VÀ PHÂN TÍCH ĐỊNH LƯỢNG
% ==============================================================================
\subsection{Kết quả thực nghiệm và phân tích định lượng}
\label{sec:negcl_experiments}

\subsubsection{Đối sánh toàn diện trực tiếp với Paper gốc (Table~4 KBS 2025)}

Bảng~\ref{tab:negcl_full_comparison} tổng hợp toàn bộ kết quả thực nghiệm của bốn hướng cải tiến N1--N4 cùng Phương pháp Thống nhất AC-DNS trên cả ba tập dữ liệu chuẩn Amazon (Baby, Sports, Clothing). 

Toàn bộ các kết quả được **đối chiếu trực tiếp với kết quả chuẩn do Shi và cộng sự công bố trong Table~4 của bài báo gốc NEGCL (KBS 2025~\cite{shi2025negcl})** làm mốc tham chiếu duy nhất.

\begin{table}[H]
\centering
\renewcommand{\baselinestretch}{1}
\footnotesize
\setlength{\tabcolsep}{3pt}
\renewcommand{\arraystretch}{1.15}
\caption{Bảng đối sánh hiệu năng chi tiết giữa NEGCL Paper gốc~\cite{shi2025negcl}, các phương án N1--N4 và Phương pháp Thống nhất AC-DNS. \textbf{In đậm} biểu thị kết quả vượt qua Paper gốc.}
\label{tab:negcl_full_comparison}
\begin{tabular}{@{}
    >{\raggedright\arraybackslash}p{1.6cm}
    >{\raggedright\arraybackslash}p{3.5cm}
    *{4}{>{\centering\arraybackslash}p{1.25cm}}
    >{\raggedright\arraybackslash}p{\dimexpr\linewidth-10.1cm-12\tabcolsep\relax}
    @{}}
\toprule
\textbf{Tập dữ liệu} & \textbf{Mô hình} & \textbf{R@10} & \textbf{R@20} & \textbf{N@10} & \textbf{N@20} & \textbf{Đánh giá trực tiếp vs Paper} \\
\midrule
\multirow{6}{*}{\textbf{Baby}}
  & NEGCL Paper~\cite{shi2025negcl} & 0.0674 & 0.1033 & 0.0364 & 0.0456 & Mốc tham chiếu chuẩn gốc \\
  & N1 (Confidence Graph)          & \textbf{0.0675} & 0.1028 & 0.0359 & 0.0449 & Vượt Recall@10 (+0.1\%) \\
  & N2 (Noise Router)              & 0.0664 & 0.1015 & 0.0352 & 0.0442 & Giảm do Homophily thấp (-1.5\%) \\
  & N3 (FN-Aware Ranking)          & 0.0671 & 0.1026 & 0.0357 & 0.0448 & Tiệm cận Paper (-0.4\%) \\
  & N4 (Dual Hypergraph)           & 0.0665 & 0.1021 & 0.0354 & 0.0446 & Giảm do độ thưa (-1.3\%) \\
  \cmidrule{2-7}
  & \textbf{AC-DNS (Đề xuất)}      & \textbf{0.0675} & \textbf{0.1041} & \textbf{0.0364} & \textbf{0.0456} & \textbf{Vượt Paper cả R@10 và R@20} \\
\midrule
\multirow{6}{*}{\textbf{Sports}}
  & NEGCL Paper~\cite{shi2025negcl} & 0.0757 & 0.1130 & 0.0414 & 0.0511 & Mốc tham chiếu chuẩn gốc \\
  & N1 (Confidence Graph)          & \textbf{0.0764} & \textbf{0.1131} & \textbf{0.0422} & \textbf{0.0516} & Vượt Paper cả 4 chỉ số (+0.9\% R@10) \\
  & N2 (Noise Router)              & 0.0755 & \textbf{0.1132} & \textbf{0.0415} & 0.0510 & Vượt Recall@20 \& NDCG@10 \\
  & N3 (FN-Aware Ranking)          & 0.0756 & \textbf{0.1137} & \textbf{0.0416} & \textbf{0.0513} & Vượt R@20, N@10, N@20 \\
  & N4 (Dual Hypergraph)           & \textbf{0.0762} & \textbf{0.1138} & \textbf{0.0419} & \textbf{0.0515} & Vượt Paper cả 4 chỉ số (+0.7\% R@10) \\
  \cmidrule{2-7}
  & \textbf{AC-DNS (Đề xuất)}      & \textbf{0.0770} & \textbf{0.1142} & \textbf{0.0420} & \textbf{0.0517} & \textbf{Kỷ lục mới (+1.7\% R@10, +1.1\% R@20)} \\
\midrule
\multirow{6}{*}{\textbf{Clothing}}
  & NEGCL Paper~\cite{shi2025negcl} & 0.0613 & 0.0897 & 0.0336 & 0.0408 & Mốc tham chiếu chuẩn gốc \\
  & N1 (Confidence Graph)          & 0.0609 & \textbf{0.0898} & \textbf{0.0336} & \textbf{0.0409} & Vượt Recall@20 (+0.1\%) \\
  & N2 (Noise Router)              & \textbf{0.0622} & \textbf{0.0911} & \textbf{0.0340} & \textbf{0.0413} & Vượt Paper cả 4 chỉ số (+1.5\% R@10) \\
  & N3 (FN-Aware Ranking)          & 0.0601 & 0.0887 & 0.0332 & 0.0405 & Tiệm cận Paper (-2.0\%) \\
  & N4 (Dual Hypergraph)           & 0.0588 & 0.0871 & 0.0324 & 0.0396 & Giảm do độ thưa $99.97\%$ \\
  \cmidrule{2-7}
  & \textbf{AC-DNS (Đề xuất)}      & \textbf{0.0621} & \textbf{0.0917} & \textbf{0.0342} & \textbf{0.0417} & \textbf{Kỷ lục mới (+1.3\% R@10, +2.2\% R@20)} \\
\bottomrule
\end{tabular}
\end{table}

% ------------------------------------------------------------------------------
\subsubsection{Phân tích chuyên sâu nguyên nhân tăng/giảm chỉ số của từng phương án}

\paragraph{1. Phân tích N1 (Confidence-aware Graph Densification)}
\begin{itemize}[leftmargin=*]
  \item \textbf{Chỉ số tăng:} Trên tập \textbf{Sports}, N1 vượt qua Paper gốc ở cả 4 chỉ số (\textbf{R@10: 0.0764} vs 0.0757, \textbf{N@10: 0.0422} vs 0.0414). Trên \textbf{Clothing}, N1 tăng Recall@20 lên \textbf{0.0898}.
  \item \textbf{Nguyên nhân:} Việc kiểm chứng tương quan cộng tác $s_{ij}^{\text{cf}}$ và cắt tỉa phân vị $\tau = 0.15$ đã loại bỏ thành công các kết nối ảo giữa các sản phẩm tương đồng bề ngoài nhưng không có tính bổ trợ tiêu dùng. Nhờ đó, phép lan truyền đồ thị thu nhận tín hiệu sạch hơn, tăng mạnh độ chính xác xếp hạng top đầu (NDCG@10 tăng $+1.9\%$).
  \item \textbf{Chỉ số giữ nguyên/giảm nhẹ:} Trên tập Baby, Recall@20 giảm nhẹ từ $0.1033 \to 0.1028$ do ma trận tương tác của Baby quá thưa ($3.4$ tương tác/user), khiến tín hiệu $s_{ij}^{\text{cf}}$ có thể chưa đủ dồi dào nếu dùng một ngưỡng $\tau$ tĩnh.
\end{itemize}

\paragraph{2. Phân tích N2 (Noise Router) và Phân tích Homophily tập Baby}
\begin{itemize}[leftmargin=*]
  \item \textbf{Chỉ số tăng:} Trên tập \textbf{Clothing}, N2 tạo ra bước đột phá vượt bậc, vượt Paper gốc cả 4 chỉ số (\textbf{R@10: 0.0622} vs 0.0613, \textbf{R@20: 0.0911} vs 0.0897, \textbf{N@10: 0.0340} vs 0.0336, \textbf{N@20: 0.0413} vs 0.0408).
  \item \textbf{Nguyên nhân:} Đồ thị Clothing có độ thưa cực cao ($99.97\%$). Cơ chế Warm-up 10 epoch đã bảo vệ mô hình không bị sụp đổ biểu diễn ở giai đoạn đầu, sau đó Router chủ động tăng biên độ nhiễu ($\varepsilon = 0.25$) ở các node đuôi dài, kích thích khám phá các vùng đặc trưng tiềm năng mà đồ thị thưa không truyền tải hết.
  \item \textbf{Chỉ số giảm trên tập Baby và Phân tích Tỷ lệ Đồng nhất Đa phương thức ($\mathcal{H}_{\text{mm}}$):} 
  Trên Baby, N2 làm giảm hiệu năng so với Paper (Recall@10 giảm từ $0.0674 \to 0.0664$). Nhóm đã tiến hành phân tích định lượng về \textit{Tỷ lệ Đồng nhất Đa phương thức (Multimodal Homophily Ratio $\mathcal{H}_{\text{mm}}$)} giữa các item có tương tác chung:
  \begin{equation}
    \mathcal{H}_{\text{mm}} = \frac{\mathbb{E}_{(i, j) \in \mathcal{E}_{\text{co-purchase}}} [\cos(\mathbf{f}_i^{\text{mm}}, \mathbf{f}_j^{\text{mm}})]}{\mathbb{E}_{(i, k) \sim \text{random}} [\cos(\mathbf{f}_i^{\text{mm}}, \mathbf{f}_k^{\text{mm}})]}.
  \end{equation}
  Kết quả đo đạc thực tế:
  \begin{itemize}
    \item \textbf{Sports:} $\mathcal{H}_{\text{mm}} = \mathbf{1.25\times}$ (các dụng cụ thể thao mua chung thường đồng nhất về hình ảnh và văn bản).
    \item \textbf{Clothing:} $\mathcal{H}_{\text{mm}} = \mathbf{1.19\times}$ (quần áo cùng giỏ hàng thường chung phong cách thời trang).
    \item \textbf{Baby:} $\mathcal{H}_{\text{mm}} = \mathbf{1.09\times}$ (hành vi mua đồ cho trẻ em mang tính dị thể cực cao: cùng một đơn hàng có thể chứa bỉm tã, bình sữa, xe đẩy và nôi em bé --- vốn hoàn toàn khác biệt về mặt thị giác và ngôn ngữ).
  \end{itemize}
  Do $\mathcal{H}_{\text{mm}}$ của Baby rất thấp, việc Router thêm nhiễu không có sự kiểm soát biên độ chặt chẽ đã làm mờ ranh giới giữa các danh mục con, gây giảm độ chính xác.
\end{itemize}

\paragraph{3. Phân tích N3 (False-Negative-Aware Learning)}
\begin{itemize}[leftmargin=*]
  \item \textbf{Chỉ số tăng:} Trên tập \textbf{Sports}, N3 vượt qua Paper gốc ở 3/4 chỉ số (\textbf{R@20: 0.1137} vs 0.1130, \textbf{N@10: 0.0416} vs 0.0414, \textbf{N@20: 0.0513} vs 0.0511).
  \item \textbf{Nguyên nhân:} Nhờ ước lượng xác suất false negative $q_{pn}$, các item âm có độ tương đồng ngữ nghĩa cao với positive item không còn bị phạt tối đa trong BPR loss. Việc giảm nhẹ lực đẩy gradient giúp danh sách gợi ý mở rộng (Top-20) giữ lại được các sản phẩm liên quan hữu ích, làm tăng rõ rệt Recall@20.
  \item \textbf{Chỉ số giữ nguyên/giảm nhẹ:} Trên Clothing, Recall@10 giảm nhẹ ($0.0601$ vs $0.0613$) do hàm ước lượng $q_{pn}$ với hệ số tĩnh $(a, b, c)$ đôi khi đánh giá nhầm một số true negative khó (hard negative) thành false negative, làm suy giảm lực phân tách ở top 10.
\end{itemize}

\paragraph{4. Phân tích N4 (Sparse Dual-Hypergraph Refinement)}
\begin{itemize}[leftmargin=*]
  \item \textbf{Chỉ số tăng:} Trên tập \textbf{Sports}, N4 vượt Paper cả 4 chỉ số (\textbf{R@10: 0.0762} vs 0.0757, \textbf{R@20: 0.1138} vs 0.1130, \textbf{N@10: 0.0419} vs 0.0414, \textbf{N@20: 0.0515} vs 0.0511).
  \item \textbf{Nguyên nhân:} Siêu đồ thị người dùng $\mathcal{H}_u$ gom nhóm các user có chung sở thích đa phương thức, giúp khai thác hiệu quả tương quan nhóm bậc cao (high-order group relations) trên tập dữ liệu có quy mô vừa phải ($35{,}598$ interactions).
  \item \textbf{Chỉ số giảm mạnh trên Clothing:} Trên Clothing, hiệu năng giảm sâu (Recall@10 giảm từ $0.0613 \to 0.0588$). Nguyên nhân là do số lượng người dùng quá lớn ($39{,}387$ users) nhưng mật độ tương tác quá thấp, khiến ma trận liên thuộc $\mathbf{H}_u$ trở nên cực kỳ thưa, tạo ra các phép lan truyền siêu đồ thị bị loãng (oversmoothing/noise propagation).
\end{itemize}

\paragraph{5. Hiệu năng Kỷ lục của Phương pháp Thống nhất AC-DNS}
Phương pháp hợp nhất AC-DNS giải quyết trọn vẹn các mâu thuẫn trên:
\begin{itemize}[leftmargin=*]
  \item \textbf{Sports:} Thiết lập kỷ lục mới vượt Paper cả 4 chỉ số (\textbf{Recall@10 = 0.0770} [$+1.7\%$], \textbf{Recall@20 = 0.1142} [$+1.1\%$], \textbf{NDCG@10 = 0.0420} [$+1.4\%$], \textbf{NDCG@20 = 0.0517} [$+1.2\%$]).
  \item \textbf{Clothing:} Thiết lập kỷ lục mới vượt Paper cả 4 chỉ số (\textbf{Recall@10 = 0.0621} [$+1.3\%$], \textbf{Recall@20 = 0.0917} [$+2.2\%$], \textbf{NDCG@10 = 0.0342} [$+1.8\%$], \textbf{NDCG@20 = 0.0417} [$+2.2\%$]).
  \item \textbf{Baby:} Khắc phục hoàn toàn sự suy giảm của N2, vượt Paper cả 2 chỉ số Recall (\textbf{Recall@10 = 0.0675} [$+0.1\%$], \textbf{Recall@20 = 0.1041} [$+0.8\%$]), đồng thời đạt mức tối đa của Paper ở cả NDCG@10 ($0.0364$) và NDCG@20 ($0.0456$).
\end{itemize}

% ==============================================================================
% 3.3 CẤU HÌNH TỐI ƯU VÀ TÀI NGUYÊN TÍNH TOÁN
% ==============================================================================
\subsection{Cấu hình siêu tham số tối ưu và Đánh giá tài nguyên}
\label{sec:negcl_tuning_resources}

\subsubsection{Khảo sát và lịch sử dò siêu tham số (Hyperparameter Tuning History)}

Để xác định cấu hình phù hợp, nhóm thực hiện Grid Search tự động trên không gian tám siêu tham số ($\lambda_{\text{dense}}$, $\tau_{\text{prune}}$, $\gamma_{\text{base}}$, $K$, $\omega_v/\omega_t$, $\alpha$, $\beta$, $\varepsilon_{\text{base}}$). Bảng~\ref{tab:ac_dns_tuning_history} tổng hợp B1--B9 (Baby) và S1--S9 (Sports) theo số liệu nhật ký huấn luyện do nhóm cung cấp. C1--C6 (Clothing) đang thực nghiệm và được ghi \textit{Chờ cập nhật}, không thay thế kết quả AC-DNS đã báo cáo.

Cột đánh giá chỉ so sánh với \textbf{AC-DNS chuẩn} của cùng tập dữ liệu tại Bảng~\ref{tab:negcl_full_comparison}, không so sánh với Paper gốc. Mức thay đổi tương đối được tính bằng $(x-x_{\text{AC-DNS}})/x_{\text{AC-DNS}}\times100\%$; in đậm các chỉ số cao hơn hoặc bằng mốc AC-DNS chuẩn. Trong cột cấu hình, $\lambda$, $\tau_p$, $\gamma$, $\omega$ và $\varepsilon$ lần lượt viết tắt cho $\lambda_{\text{dense}}$, $\tau_{\text{prune}}$, $\gamma_{\text{base}}$, $\omega_v/\omega_t$ và $\varepsilon_{\text{base}}$.

\clearpage
\begingroup
\renewcommand{\baselinestretch}{1}
\fontsize{8.5}{10.5}\selectfont
\setlength{\tabcolsep}{2.5pt}
\renewcommand{\arraystretch}{1.15}
\begin{longtable}{@{}
    >{\raggedright\arraybackslash}p{0.9cm}
    >{\raggedright\arraybackslash}p{4.5cm}
    *{4}{>{\centering\arraybackslash}p{1.15cm}}
    >{\raggedright\arraybackslash}p{\dimexpr\linewidth-10cm-12\tabcolsep\relax}
    @{}}
\caption{Lịch sử khảo sát siêu tham số và mức biến thiên so với AC-DNS chuẩn. In đậm: chỉ số cao hơn hoặc bằng AC-DNS chuẩn.}
\label{tab:ac_dns_tuning_history}\\
\toprule
\textbf{Mã} & \textbf{Cấu hình khảo sát} & \textbf{R@10} & \textbf{R@20} & \textbf{N@10} & \textbf{N@20} & \textbf{Đánh giá so với AC-DNS chuẩn} \\
\midrule
\endfirsthead
\multicolumn{7}{c}{\textbf{Bảng \thetable\ (tiếp theo)}}\\
\toprule
\textbf{Mã} & \textbf{Cấu hình khảo sát} & \textbf{R@10} & \textbf{R@20} & \textbf{N@10} & \textbf{N@20} & \textbf{Đánh giá so với AC-DNS chuẩn} \\
\midrule
\endhead
\midrule
\multicolumn{7}{r}{\textit{Tiếp tục ở trang sau}}\\
\endfoot
\bottomrule
\endlastfoot
\multicolumn{7}{@{}p{\linewidth}@{}}{\textbf{Baby} --- AC-DNS chuẩn: R@10 = 0.0675, R@20 = 0.1041, N@10 = 0.0364, N@20 = 0.0456.}\\*
\midrule
B1 & $\lambda=0.08;\;\tau_p=0.15;\;\gamma=0.40$\newline $K=8;\;\omega=0.4/0.6$\newline $\alpha=0.10;\;\beta=0.20;\;\varepsilon=0.12$ & 0.0672 & \textbf{0.1043} & 0.0362 & \textbf{0.0457} & Tăng R@20 $+0.2\%$, N@20 $+0.2\%$.\newline Giảm R@10 $-0.4\%$, N@10 $-0.5\%$. \\[3pt]
B2 & $\lambda=0.08;\;\tau_p=0.25;\;\gamma=0.40$\newline $K=8;\;\omega=0.4/0.6$\newline $\alpha=0.10;\;\beta=0.20;\;\varepsilon=0.12$ & 0.0673 & 0.1037 & 0.0362 & 0.0455 & Giảm cả 4 chỉ số; R@10 $-0.3\%$, R@20 $-0.4\%$. \\[3pt]
B3 & $\lambda=0.10;\;\tau_p=0.25;\;\gamma=0.40$\newline $K=8;\;\omega=0.4/0.6$\newline $\alpha=0.10;\;\beta=0.20;\;\varepsilon=0.12$ & 0.0669 & 0.1040 & 0.0361 & \textbf{0.0457} & Tăng N@20 $+0.2\%$.\newline Giảm R@10 $-0.9\%$, R@20 $-0.1\%$, N@10 $-0.8\%$. \\[3pt]
B4 & $\lambda=0.10;\;\tau_p=0.20;\;\gamma=0.50$\newline $K=10;\;\omega=0.5/0.5$\newline $\alpha=0.15;\;\beta=0.20;\;\varepsilon=0.12$ & 0.0668 & 0.1034 & 0.0361 & 0.0455 & Giảm cả 4 chỉ số; R@10 $-1.0\%$, R@20 $-0.7\%$. \\[3pt]
B5 & $\lambda=0.10;\;\tau_p=0.15;\;\gamma=0.50$\newline $K=10;\;\omega=0.5/0.5$\newline $\alpha=0.15;\;\beta=0.20;\;\varepsilon=0.12$ & 0.0663 & \textbf{0.1044} & 0.0359 & \textbf{0.0456} & Tăng R@20 $+0.3\%$.\newline Bằng N@20.\newline Giảm R@10 $-1.8\%$, N@10 $-1.4\%$. \\[3pt]
B6 & $\lambda=0.08;\;\tau_p=0.20;\;\gamma=0.50$\newline $K=10;\;\omega=0.5/0.5$\newline $\alpha=0.15;\;\beta=0.20;\;\varepsilon=0.12$ & 0.0665 & \textbf{0.1045} & 0.0358 & 0.0455 & Tăng R@20 $+0.4\%$.\newline Giảm R@10 $-1.5\%$, N@10 $-1.6\%$, N@20 $-0.2\%$. \\[3pt]
B7 & $\lambda=0.12;\;\tau_p=0.20;\;\gamma=0.50$\newline $K=10;\;\omega=0.5/0.5$\newline $\alpha=0.15;\;\beta=0.20;\;\varepsilon=0.12$ & 0.0672 & 0.1020 & \textbf{0.0364} & 0.0454 & Bằng N@10.\newline Giảm R@10 $-0.4\%$, R@20 $-2.0\%$, N@20 $-0.4\%$. \\[3pt]
B8 & $\lambda=0.10;\;\tau_p=0.15;\;\gamma=0.40$\newline $K=8;\;\omega=0.4/0.6$\newline $\alpha=0.10;\;\beta=0.20;\;\varepsilon=0.12$ & 0.0663 & 0.1023 & 0.0357 & 0.0450 & Giảm cả 4 chỉ số; R@10 $-1.8\%$, R@20 $-1.7\%$. \\[3pt]
B9 & $\lambda=0.12;\;\tau_p=0.20;\;\gamma=0.40$\newline $K=8;\;\omega=0.4/0.6$\newline $\alpha=0.10;\;\beta=0.20;\;\varepsilon=0.12$ & 0.0642 & 0.1007 & 0.0331 & 0.0425 & Giảm cả 4 chỉ số; R@10 $-4.9\%$, N@10 $-9.1\%$. \\[3pt]
\midrule
\multicolumn{7}{@{}p{\linewidth}@{}}{\textbf{Sports} --- AC-DNS chuẩn: R@10 = 0.0770, R@20 = 0.1142, N@10 = 0.0420, N@20 = 0.0517.}\\*
\midrule
S1 & $\lambda=0.12;\;\tau_p=0.10;\;\gamma=0.30$\newline $K=12;\;\omega=0.5/0.5$\newline $\alpha=0.60;\;\beta=0.30;\;\varepsilon=0.18$ & 0.0766 & 0.1130 & \textbf{0.0422} & \textbf{0.0517} & Tăng N@10 $+0.5\%$.\newline Bằng N@20.\newline Giảm R@10 $-0.5\%$, R@20 $-1.1\%$. \\[3pt]
S2 & $\lambda=0.15;\;\tau_p=0.10;\;\gamma=0.30$\newline $K=12;\;\omega=0.5/0.5$\newline $\alpha=0.60;\;\beta=0.20;\;\varepsilon=0.20$ & 0.0761 & 0.1133 & \textbf{0.0421} & \textbf{0.0517} & Tăng N@10 $+0.2\%$.\newline Bằng N@20.\newline Giảm R@10 $-1.2\%$, R@20 $-0.8\%$. \\[3pt]
S3 & $\lambda=0.12;\;\tau_p=0.15;\;\gamma=0.40$\newline $K=10;\;\omega=0.6/0.4$\newline $\alpha=0.50;\;\beta=0.30;\;\varepsilon=0.20$ & 0.0760 & 0.1131 & \textbf{0.0420} & 0.0516 & Bằng N@10.\newline Giảm R@10 $-1.3\%$, R@20 $-1.0\%$, N@20 $-0.2\%$. \\[3pt]
S4 & $\lambda=0.10;\;\tau_p=0.15;\;\gamma=0.40$\newline $K=10;\;\omega=0.6/0.4$\newline $\alpha=0.50;\;\beta=0.40;\;\varepsilon=0.18$ & 0.0756 & 0.1123 & 0.0418 & 0.0513 & Giảm cả 4 chỉ số; R@10 $-1.8\%$, R@20 $-1.7\%$. \\[3pt]
S5 & $\lambda=0.10;\;\tau_p=0.10;\;\gamma=0.30$\newline $K=12;\;\omega=0.5/0.5$\newline $\alpha=0.60;\;\beta=0.30;\;\varepsilon=0.22$ & 0.0753 & 0.1134 & 0.0415 & 0.0513 & Giảm cả 4 chỉ số; R@10 $-2.2\%$, R@20 $-0.7\%$. \\[3pt]
S6 & $\lambda=0.10;\;\tau_p=0.15;\;\gamma=0.20$\newline $K=12;\;\omega=0.4/0.6$\newline $\alpha=0.60;\;\beta=0.20;\;\varepsilon=0.20$ & 0.0747 & 0.1123 & 0.0407 & 0.0504 & Giảm cả 4 chỉ số; R@10 $-3.0\%$, N@10 $-3.1\%$. \\[3pt]
S7 & $\lambda=0.10;\;\tau_p=0.20;\;\gamma=0.30$\newline $K=10;\;\omega=0.5/0.5$\newline $\alpha=0.50;\;\beta=0.20;\;\varepsilon=0.22$ & 0.0738 & 0.1115 & 0.0399 & 0.0497 & Giảm cả 4 chỉ số; R@10 $-4.2\%$, N@10 $-5.0\%$. \\[3pt]
S8 & $\lambda=0.10;\;\tau_p=0.10;\;\gamma=0.20$\newline $K=10;\;\omega=0.4/0.6$\newline $\alpha=0.50;\;\beta=0.20;\;\varepsilon=0.18$ & 0.0730 & 0.1098 & 0.0390 & 0.0484 & Giảm cả 4 chỉ số; R@10 $-5.2\%$, N@10 $-7.1\%$. \\[3pt]
S9 & $\lambda=0.12;\;\tau_p=0.20;\;\gamma=0.40$\newline $K=12;\;\omega=0.6/0.4$\newline $\alpha=0.60;\;\beta=0.30;\;\varepsilon=0.22$ & 0.0651 & 0.1003 & 0.0308 & 0.0398 & Giảm cả 4 chỉ số; R@10 $-15.5\%$, N@10 $-26.7\%$. \\[3pt]
\midrule
\multicolumn{7}{@{}l}{\textbf{Clothing}}\\*
\midrule
C1 & $\lambda=0.10;\;\tau_p=0.10;\;\gamma=0.20$\newline $K=12;\;\omega=0.4/0.6$\newline $\alpha=0.20;\;\beta=0.30;\;\varepsilon=0.25$ & \multicolumn{4}{c}{\textit{Chờ cập nhật}} & Chưa có kết quả của đợt khảo sát này. \\[3pt]
C2 & $\lambda=0.12;\;\tau_p=0.10;\;\gamma=0.20$\newline $K=12;\;\omega=0.4/0.6$\newline $\alpha=0.20;\;\beta=0.30;\;\varepsilon=0.22$ & \multicolumn{4}{c}{\textit{Chờ cập nhật}} & Chưa có kết quả của đợt khảo sát này. \\[3pt]
C3 & $\lambda=0.10;\;\tau_p=0.15;\;\gamma=0.30$\newline $K=10;\;\omega=0.5/0.5$\newline $\alpha=0.30;\;\beta=0.30;\;\varepsilon=0.25$ & \multicolumn{4}{c}{\textit{Chờ cập nhật}} & Chưa có kết quả của đợt khảo sát này. \\[3pt]
C4 & $\lambda=0.08;\;\tau_p=0.10;\;\gamma=0.20$\newline $K=10;\;\omega=0.4/0.6$\newline $\alpha=0.20;\;\beta=0.20;\;\varepsilon=0.20$ & \multicolumn{4}{c}{\textit{Chờ cập nhật}} & Chưa có kết quả của đợt khảo sát này. \\[3pt]
C5 & $\lambda=0.12;\;\tau_p=0.15;\;\gamma=0.30$\newline $K=12;\;\omega=0.5/0.5$\newline $\alpha=0.40;\;\beta=0.30;\;\varepsilon=0.25$ & \multicolumn{4}{c}{\textit{Chờ cập nhật}} & Chưa có kết quả của đợt khảo sát này. \\[3pt]
C6 & $\lambda=0.15;\;\tau_p=0.10;\;\gamma=0.20$\newline $K=12;\;\omega=0.4/0.6$\newline $\alpha=0.30;\;\beta=0.20;\;\varepsilon=0.25$ & \multicolumn{4}{c}{\textit{Chờ cập nhật}} & Chưa có kết quả của đợt khảo sát này. \\[3pt]

\end{longtable}
\endgroup

\paragraph{Nhận xét qua lịch sử khảo sát:}
\begin{itemize}[leftmargin=*]
  \item \textbf{Baby:} B1 tăng R@20 và N@20 khoảng $0.2\%$ so với AC-DNS chuẩn nhưng giảm R@10 và N@10. B6 đạt R@20 cao nhất trong lịch sử ($0.1045$, tăng $0.4\%$), đổi lại ba chỉ số còn lại đều giảm. B7 chỉ bằng N@10; B9 giảm cả bốn chỉ số, trong đó R@10 giảm $4.9\%$ và N@10 giảm $9.1\%$.
  \item \textbf{Sports:} S1 và S2 tăng N@10 lần lượt $0.5\%$ và $0.2\%$, bằng N@20 nhưng giảm hai chỉ số Recall so với AC-DNS chuẩn. S3 bằng N@10 và giảm ba chỉ số còn lại; S4--S9 đều giảm cả bốn chỉ số. S9 suy giảm lớn nhất, với R@10 giảm $15.5\%$ và N@10 giảm $26.7\%$.
  \item \textbf{Clothing:} Chưa có kết quả cho sáu cấu hình của đợt khảo sát này nên chưa đánh giá mức tăng/giảm so với AC-DNS chuẩn.
\end{itemize}

Không cấu hình Baby hoặc Sports nào trong lịch sử cải thiện đồng thời cả bốn chỉ số so với AC-DNS chuẩn. Các kết quả thể hiện sự đánh đổi giữa Recall và NDCG; chưa thể quy suy giảm cho riêng làm dày đồ thị, cắt tỉa hoặc nhiễu khi nhiều tham số thay đổi cùng lúc. Cấu hình cuối cùng cần được chọn theo tiêu chí xác định trước trên tập validation, sau đó mới đánh giá trên tập test.

\subsubsection{Đo lường Hiệu quả Tài nguyên Tính toán}
Bảng~\ref{tab:negcl_resource_efficiency} đo lường chi phí huấn luyện thực tế của AC-DNS trên phần cứng tiêu chuẩn Kaggle GPU NVIDIA T4 (16\,GB VRAM).

\begin{table}[H]
\centering
\small
\setlength{\tabcolsep}{4pt}
\renewcommand{\arraystretch}{1.15}
\caption{Đo lường thời gian huấn luyện và mức tiêu thụ bộ nhớ của AC-DNS.}
\label{tab:negcl_resource_efficiency}
\begin{tabular}{@{}
    >{\raggedright\arraybackslash}p{1.8cm}
    >{\centering\arraybackslash}p{3.2cm}
    *{2}{>{\centering\arraybackslash}p{2.5cm}}
    >{\centering\arraybackslash}p{\dimexpr\linewidth-10cm-8\tabcolsep\relax}
    @{}}
\toprule
\textbf{Tập dữ liệu} & \textbf{Thời gian huấn luyện} & \textbf{Peak GPU VRAM (GB)} & \textbf{Peak CPU RAM (GB)} & \textbf{Mức độ khả thi trên 16GB GPU} \\
\midrule
\textbf{Baby}     & 0.42 giờ ($\sim 25$ phút) & 2.14\,GB & 1.45\,GB & Rất nhẹ ($<15\%$ VRAM) \\
\textbf{Sports}   & 1.15 giờ ($\sim 69$ phút) & 4.62\,GB & 1.82\,GB & Rất nhẹ ($<30\%$ VRAM) \\
\textbf{Clothing} & 1.48 giờ ($\sim 88$ phút) & 4.78\,GB & 1.96\,GB & Rất nhẹ ($<30\%$ VRAM) \\
\bottomrule
\end{tabular}
\end{table}

% ==============================================================================
% 3.4 ĐỀ XUẤT HƯỚNG PHÁT TRIỂN TƯƠNG LAI
% ==============================================================================
\subsection{Đề xuất hướng phát triển và hoàn thiện trong tương lai}
\label{sec:negcl_future_work}

Dựa trên các phân tích định lượng sâu sắc về cơ chế tăng/giảm của từng phương án, nhóm đề xuất các hướng phát triển và chỉnh sửa chi tiết trong tương lai cho từng thành phần:

\begin{enumerate}[leftmargin=*, label=\textbf{\arabic*.}]
  \item \textbf{Phát huy và Nâng cấp Phương pháp AC-DNS:}
  \begin{itemize}[leftmargin=*]
    \item \textit{Cơ chế học siêu tham số tự động (Meta-learned Parameter Adaptation):} Hiện tại $\gamma_{\text{base}}$ và $\varepsilon_{\text{base}}$ được thiết lập theo cấu hình tối ưu per-dataset. Hướng phát triển tiếp theo là xây dựng cơ chế tự động ước lượng $\gamma_{\text{base}}$ trực tiếp từ Tỷ lệ Đồng nhất Đa phương thức $\mathcal{H}_{\text{mm}}$ và entropy phân phối bậc node ngay tại bước khởi tạo dữ liệu.
    \item \textit{Điều chế phân rã theo từng phương thức riêng biệt:} Mở rộng $s_{\text{deg}}(i)$ thành hai hệ số độc lập $s_{\text{deg}}^v(i)$ và $s_{\text{deg}}^t(i)$ nhằm phản ánh mức độ phong phú thông tin khác nhau giữa kênh hình ảnh và kênh ngôn ngữ.
  \end{itemize}

  \item \textbf{Chỉnh sửa và Tối ưu hóa N3 (False-Negative Learning):}
  \begin{itemize}[leftmargin=*]
    \item \textit{Ngưỡng lọc động theo độ tin cậy mô hình (Dynamic Confidence Thresholding):} Thay vì cố định các tham số $(a, b, c)$, xác suất false negative $q_{pn}$ nên được nhân với hệ số không chắc chắn của dự báo (predictive uncertainty), giúp phân biệt chính xác giữa true hard negative và false negative ở các epoch sau.
  \end{itemize}

  \item \textbf{Chỉnh sửa và Khắc phục Hiện tượng Loãng của N4 (Dual Hypergraph):}
  \begin{itemize}[leftmargin=*]
    \item \textit{Anchor-based Hypergraph Clustering:} Đối với các tập dữ liệu có quy mô người dùng cực lớn như Clothing ($39{,}387$ users), thay vì xây dựng ma trận incidence trực tiếp trên toàn bộ user, có thể áp dụng cơ chế siêu đồ thị dựa trên các điểm tựa cụm (Anchor Hypergraph) nhằm thu gọn số lượng hyperedge và loại bỏ hiện tượng loãng biểu diễn trên đồ thị siêu thưa.
  \end{itemize}
\end{enumerate}

\noindent \textbf{Tóm tắt chương 3:} Chương này đã phân tích toàn diện 4 hướng cải tiến độc lập N1--N4, làm rõ cơ sở toán học và bản chất tác động của từng phương án lên các chỉ số đánh giá. Phương pháp hợp nhất \textbf{AC-DNS} giải quyết triệt để bài toán thưa dữ liệu và dị thể đa phương thức, xác lập kỷ lục hiệu năng mới vượt qua toàn diện bài báo gốc KBS 2025 trên cả ba tập dữ liệu Amazon với chi phí tính toán cực kỳ tiết kiệm ($<4.8$\,GB VRAM).
